\section{Q2 利润桥：预告已确认的部分}

公司《2026 年半年度业绩预告》直接披露 Q2 归母净利润区间，而非仅由研究者倒算。以 Q1 已披露归母净利润 3.388213 亿元对照，H1 预告与 Q2 区间相互勾稽；Q2 扣非归母净利润则是 H1 扣非预告减 Q1 已披露值的派生数字，应与公司直接披露严格区分。

\begin{exhibitbox}[Q1--Q2--H1 利润桥：只做披露项之间的勾稽]
\begin{tabularx}{\exhibitboxwidth}{L{3.1cm}R{2.4cm}R{2.7cm}R{2.7cm}X}
\toprule
项目 & 2026Q1 已披露 & 2026H1 预告 & 2026Q2 区间 & 口径 \\
\midrule
归母净利润 & 3.3882 亿元 & 11.00--13.00 亿元 & 7.6118--9.6118 亿元 & Q2 区间由公司直接披露 \\
扣非归母净利润 & 3.5656 亿元 & 11.25--13.15 亿元 & 7.6844--9.5844 亿元 & Q2 为 H1 减 Q1 的派生值 \\
基本每股收益 & 0.2964 元 & 0.9621--1.1370 元 & 不单独推导 & 加权股本不同，不作二次建模 \\
Q2 环比归母 & -- & -- & +124.65\%--+183.68\% & 公司公告直接披露 \\
\bottomrule
\end{tabularx}
\end{exhibitbox}
\sourcenote{雅化集团《2026 年第一季度报告》（2026-04-27）第 2 页；《2026 年半年度业绩预告》（2026-07-07）第 1 页。}

这张桥只回答“利润是否跳升”，不回答“为什么跳升了多少”。公司给出的驱动是锂盐市场价格、销量和销售均价同步增长，以及运营效率、成本控制和矿--产--销平衡；H1 预告没有披露收入、毛利、套保、分部利润和现金流。因此，本研究把 Q2 归母区间作为已确认事件，把驱动的金额拆分保留给正式 H1。

\section{预告相对外部全年预测：H2 才是预期分歧}

三份完整券商 PDF 均早于或在 H1 预告日发布。它们不是公司指引，也不是本报告预测；它们的价值在于显示市场已经需要 H2 达到什么水平。预告 H1 的低/高端分别代入每家的全年归母预测，得到其尚未实现的 H2 利润需求。

\begin{exhibitbox}[H1 预告对外部全年预测的覆盖与隐含 H2]
\begin{tabularx}{\exhibitboxwidth}{L{2.1cm}R{2.5cm}R{2.0cm}R{2.1cm}R{2.1cm}X}
\toprule
机构 / 报告日 & 2026E 收入 / 归母 & H1 覆盖全年 & 隐含 H2 归母 & H2 单季均值 & 事件含义 \\
\midrule
国信\\4月30日 & 161.64 / 23.15 亿元 & 47.5\%--56.2\% & 10.15--12.15 亿元 & 5.08--6.08 亿元 & H2 要求低于预告 Q2 区间 \\
开源\\4月29日 & 168.18 / 25.69 亿元 & 42.8\%--50.6\% & 12.69--14.69 亿元 & 6.35--7.35 亿元 & 需要维持较高盈利 \\
东吴\\7月7日 & 170.52 / 30.32 亿元 & 36.3\%--42.9\% & 17.32--19.32 亿元 & 8.66--9.66 亿元 & H2 接近或高于预告 Q2 \\
\bottomrule
\end{tabularx}
\end{exhibitbox}
\sourcenote{国信证券《财报点评：锂盐业务充分享受锂价弹性，民爆业务国际化布局优势明显》（2026-04-30）第 1 页；开源证券《Q1 锂盐价格延续上行，带动公司业绩高速增长》（2026-04-29）第 1 页；东吴证券《26H1 业绩预告点评：Q2 锂盐量利齐升，略超我们预期》（2026-07-07）第 1--2 页。H1 覆盖和 H2 为本研究基于公告区间的算术推导。}

\begin{keyinsight}[预期差的正确表述]
Q2 的高利润对三家外部预测都不是反证；真正的差异在 H2。国信的剩余要求相对温和，东吴的全年判断要求下半年继续维持接近 Q2 高位的单季利润。正式 H1 即使确认预告，也不能自动确认任一家全年预测；它只能通过收入、现金流、库存与单位经济披露来提高或降低这些 H2 假设的可信度。
\end{keyinsight}

\section{外部假设必须经正式披露验证}

东吴报告估计 Q2 锂盐出货 2--3 万吨、吨利润 2.6--3.4 万元、全年权益自供比例约 25\%，并估计民爆 Q2 利润约 1 亿元。上述数字对理解市场可能交易的弹性很有帮助，但均为分析师估计；本报告不给它们任何“公司已实现”身份，也不把它们填进自有 EPS 或目标价。

\begin{sourcequalitybox}[事件研究的证据层级]
一级资料确认 Q2 合并利润区间和管理层的定性归因。二级券商资料只定义待验证的 H2 假设。行业现货价、设计产能、客户名单和资源权益都不能替代雅化的 H1 销量、实际 ASP、单位成本、分部利润和经营现金流。
\end{sourcequalitybox}
